\chapter{GIỚI THIỆU}

\section{Đặt vấn đề}

Trong kỷ nguyên số hóa và cuộc cách mạng công nghiệp 4.0, Internet vạn vật đã vươn lên trở thành một trong những xu hướng công nghệ mang tính định hình tương lai. Sự bùng nổ của IoT không chỉ giới hạn ở quy mô công nghiệp mà đã len lỏi vào từng khía cạnh của đời sống con người, tiêu biểu nhất là sự phát triển mạnh mẽ của hệ sinh thái nhà thông minh. Theo các báo cáo thị trường gần đây, nhu cầu tự động hóa không gian sống đang gia tăng theo cấp số nhân, khi con người ngày càng mong muốn một môi trường sống tiện nghi, an toàn và thông minh hơn. Trong hệ sinh thái đó, hệ thống an ninh và đặc biệt là hệ thống kiểm soát cửa ra vào đóng vai trò như người gác cổng đầu tiên, bảo vệ tài sản và sự riêng tư của gia chủ. Chính vì thế, việc nghiên cứu và nâng cấp các thiết bị khóa cửa thông minh không chỉ là một nhu cầu thiết thực của thị trường mà còn là một bài toán kỹ thuật hấp dẫn, thu hút sự quan tâm của nhiều nhà nghiên cứu và kỹ sư phát triển sản phẩm.

Mặc dù đóng vai trò vô cùng quan trọng, các phương thức khóa cửa truyền thống hiện nay đang bộc lộ nhiều điểm yếu chí mạng, không còn đáp ứng đủ các tiêu chuẩn an ninh khắt khe của cuộc sống hiện đại. Phân tích một cách chuyên sâu, sự hạn chế này thể hiện rõ qua ba phương thức phổ biến nhất. Thứ nhất, đối với khóa cơ khí truyền thống, rủi ro lớn nhất nằm ở việc người dùng dễ dàng bị sao chép chìa khóa vật lý, hệ thống khóa dễ bị phá hoại bằng các công cụ cắt/cạy chuyên dụng, và sự bất tiện phiền toái khi người dùng làm mất hoặc quên mang theo chìa. Thứ hai, với các loại khóa sử dụng mật mã , nguy cơ lộ lọt thông tin là rất cao. Kẻ gian có thể dễ dàng nhìn trộm mật khẩu qua vai hoặc sử dụng các kỹ thuật phân tích dấu vết để lại trên bàn phím mờ để dò ra dãy số mở cửa. Thứ ba, đối với khóa sử dụng thẻ từ, dù mang lại sự tiện lợi hơn nhờ thao tác chạm mở nhanh chóng, nhưng bản chất thẻ từ vẫn là một vật mang vật lý. Người dùng vẫn phải đối mặt với nguy cơ làm rơi, mất cắp, hoặc bị tin tặc sử dụng các thiết bị đọc trộm để sao chép dữ liệu sóng vô tuyến từ xa một cách âm thầm mà nạn nhân không hề hay biết. Những lỗ hổng này đặt ra một yêu cầu bức thiết về một phương thức xác thực hoàn toàn mới, gắn liền với định danh duy nhất của mỗi cá nhân.

Để giải quyết triệt để những hạn chế của các phương thức mở khóa vật lý, công nghệ sinh trắc học – đặc biệt là nhận diện khuôn mặt – nổi lên như một xu hướng tất yếu nhờ khả năng xác thực tự nhiên, không cần tiếp xúc vật lý và không thể bị đánh cắp hay bỏ quên. Tuy nhiên, việc áp dụng công nghệ này vào các thiết bị IoT cỡ nhỏ lại vấp phải những rào cản kỹ thuật đáng kể. Phần lớn các hệ thống nhận diện khuôn mặt có độ chính xác cao hiện nay đều hoạt động dựa trên điện toán đám mây. Cách tiếp cận này đòi hỏi thiết bị phải liên tục gửi dữ liệu hình ảnh lên máy chủ, gây tiêu tốn băng thông mạng, tạo ra độ trễ  trong quá trình xử lý, và nghiêm trọng nhất là làm dấy lên những lo ngại sâu sắc về quyền riêng tư khi hình ảnh sinh trắc học nhạy cảm bị lưu trữ trên các nền tảng ngoại vi. Ở chiều hướng ngược lại, để đưa AI về chạy cục bộ, các nhà phát triển thường phải sử dụng đến các vi xử lý mạnh mẽ như NVIDIA Jetson Nano hay Raspberry Pi. Mặc dù mang lại hiệu năng tính toán cao, nhưng chi phí linh kiện đắt đỏ, kích thước phần cứng cồng kềnh và mức tiêu thụ điện năng lớn khiến các bo mạch này khó có thể được tích hợp vào một ổ khóa cửa nhỏ gọn, chạy bằng pin và hướng tới khả năng thương mại hóa đại chúng.

Đứng trước bài toán nan giải về sự đánh đổi giữa hiệu năng xử lý AI, chi phí phần cứng và tính bảo mật riêng tư, sự ra đời của xu hướng Edge AI và đặc biệt là TinyML đã mang đến một giải pháp "cứu cánh" mang tính đột phá. Bằng cách tối ưu hóa các mô hình mạng nơ-ron để vận hành mượt mà trên các chip vi điều khiển giới hạn tài nguyên, TinyML cho phép quá trình suy luận diễn ra hoàn toàn tại biên mạng. Tiêu biểu trong phân khúc này là vi điều khiển ESP32-S3, một dòng chip tối ưu chi phí nhưng được tích hợp tập lệnh xử lý vector dành riêng cho tính toán AI, đi kèm kết nối Wi-Fi/Bluetooth vô cùng mạnh mẽ. Nhận thấy tiềm năng to lớn từ việc kết hợp giữa công nghệ TinyML và sức mạnh của ESP32-S3 để giải bài toán an ninh nhà ở, đề tài: \textbf{"Xây dựng hệ thống nhận diện khuôn mặt cho nhà thông minh"} đã được lựa chọn. Nghiên cứu này không chỉ hướng tới việc xây dựng một sản phẩm IoT an toàn, nhận diện nhanh chóng, bảo vệ dữ liệu tuyệt đối mà còn tối ưu hóa chi phí, mở ra hướng đi mới cho các thiết bị nhà thông minh Make in Vietnam.
\section{Mục tiêu của đề tài}
Mục tiêu tổng quát của đề tài là nghiên cứu, thiết kế và chế tạo thành công một hệ thống khóa cửa thông minh tích hợp công nghệ nhận diện khuôn mặt hoạt động hoàn toàn độc lập trên nền tảng vi điều khiển ESP32-S3. Sản phẩm hướng tới việc giải quyết bài toán phụ thuộc vào điện toán đám mây, qua đó tạo ra một thiết bị an ninh cục bộ có khả năng bảo mật cao, phản hồi nhanh chóng và bảo vệ tuyệt đối quyền riêng tư sinh trắc học của người dùng. Đồng thời, giải pháp phải duy trì được mức chi phí sản xuất thấp, tạo tiền đề để dễ dàng tiếp cận và thương mại hóa trong đại đa số các hộ gia đình.

Để hiện thực hóa mục tiêu tổng quát nêu trên, đề tài xác định và cam kết hoàn thành 4 mục tiêu cụ thể và các chỉ tiêu kỹ thuật tương ứng.

Thứ nhất, về phần cứng, đề tài sẽ sử dụng vi điều khiển ESP32-S3 cùng module camera OV3660 nhằm đảm nhiệm vai trò thu thập dữ liệu hình ảnh trực tiếp tại biên. Hệ thống phần cứng phải đảm bảo khả năng tính toán điều khiển chính xác. Khối thiết bị yêu cầu phải hoạt động ổn định trong các điều kiện môi trường thực tế, có khả năng chống nhiễu điện từ.

Thứ hai, về trí tuệ nhân tạo, đề tài sẽ triển khai và làm chủ thành công các thuật toán mạng nơ-ron học sâu trên nền tảng thiết bị hạn chế tài nguyên. Đề tài cần trích xuất và tối ưu hóa mô hình nhận diện khuôn mặt để nhúng trực tiếp vào không gian bộ nhớ khiêm tốn của ESP32-S3. Tiêu chuẩn kỹ thuật đặt ra là hệ thống phải đạt độ chính xác nhận diện trên 80\% trong điều kiện ánh sáng tiêu chuẩn, đồng thời giới hạn thời gian suy luận dưới 2 giây cho mỗi lượt quét nhằm mang lại trải nghiệm mở khóa mượt mà, không có độ trễ cho người dùng.

Thứ ba, về phần mềm quản lý, đề tài sẽ xây dựng một hệ sinh thái phần mềm đồng bộ bao gồm firmware chạy trên vi điều khiển và một ứng dụng quản lý từ xa trên điện thoại thông minh và trên máy tính cá nhân. Hệ thống phải cung cấp giao diện trực quan cho phép quản trị viên thao tác các chức năng cốt lõi: đăng ký thêm khuôn mặt mới, xóa quyền truy cập, và giám sát lịch sử ra vào theo thời gian thực. Quá trình truyền tải thông tin giữa ESP32-S3 và ứng dụng phải được đảm bảo bằng các giao thức an toàn, đồng thời hỗ trợ gửi cảnh báo tự động ngay khi phát hiện nỗ lực xâm nhập bằng khuôn mặt lạ lặp lại nhiều lần.

Cuối cùng, về tối ưu hóa, đề tài sẽ tối ưu hóa toàn diện bài toán năng lượng và tài nguyên lưu trữ để đảm bảo tính ứng dụng lâu dài và bền bỉ của sản phẩm. Cụ thể, thiết bị phải được lập trình để khai thác triệt để chế độ ngủ sâu của chip ESP32-S3, tinh chỉnh các chu kỳ thức - ngủ nhằm giảm mức tiêu thụ điện năng ở trạng thái chờ xuống mức tối thiểu (cỡ $\mu$A). Song song với đó, cơ sở dữ liệu cục bộ phải được cấu trúc hợp lý để lưu trữ an toàn các vector đặc trưng khuôn mặt, đảm bảo tốc độ truy xuất nhanh và duy trì tuổi thọ cho bộ nhớ vật lý của toàn hệ thống.

\section{Đối tượng và phạm vi nghiên cứu}

\subsection{Đối tượng nghiên cứu}
Đối tượng nghiên cứu chính của đề tài là hệ sinh thái phần cứng và phần mềm cho bài toán nhận diện khuôn mặt trên thiết bị biên.

Về phần cứng, đề tài tập trung vào hai thành phần cốt lõi.

Thứ nhất là vi điều khiển ESP32-S3, một dòng chip SoC hỗ trợ Wi-Fi và Bluetooth Low Energy 5.0, được trang bị các tập lệnh mở rộng phục vụ tăng tốc thuật toán học sâu. Với vi xử lý lõi kép Xtensa 32-bit LX7 xung nhịp 240 MHz, 512 KB SRAM, 8 MB PSRAM và 16MB Flash, nó đáp ứng khả năng tính toán ma trận cần thiết cho các mô hình như MobileFaceNet.

Thành phần thứ hai là module camera OV3660, một cảm biến hình ảnh CMOS 3 Megapixel có độ nhạy sáng cao, kết nối trực tiếp với ESP32-S3 qua giao diện DVP để truyền tải dữ liệu khung hình ổn định.

Về phần mềm, đề tài nghiên cứu các thuật toán và khung phần mềm như mô hình học sâu MobileFaceNet và MTMN, được tối ưu hóa trên nền tảng ESP-DL và ESP-WHO để thực hiện bài toán nhận diện khuôn mặt thời gian thực ngay trên thiết bị.

\subsection{Phạm vi và giới hạn}
Hệ thống được thiết kế và kiểm chứng hoạt động trong môi trường văn phòng hoặc hộ gia đình với điều kiện chiếu sáng tiêu chuẩn. Phạm vi nhận diện khuôn mặt được giới hạn trong khoảng cách từ 30 đến 80 cm, đây là khoảng cách tối ưu cho các thao tác mở khóa cửa tại chỗ. Hệ thống tập trung tối đa vào kỹ thuật tính toán tại biên, do đó năng lực lưu trữ và xử lý bị giới hạn bởi phần cứng nhúng. Cụ thể, bộ nhớ Flash của ESP32-S3 được tối ưu để lưu trữ dữ liệu đặc trưng cho từ 50 đến 100 người dùng nhằm đảm bảo độ trễ nhận diện luôn dưới 2 giây. Đề tài không hướng tới việc xây dựng cơ sở dữ liệu khuôn mặt quy mô lớn hoặc các hệ thống nhận diện từ xa.

Việc truyền tin và đẩy thông báo được thực hiện thông qua hạ tầng MQTT và dịch vụ Firebase Cloud Messaging với giả định hệ thống duy trì kết nối mạng Wi-Fi ổn định. Các yếu tố về chống giả mạo cao cấp như nhận diện độ sâu bằng camera chuyên dụng hay cảm biến hồng ngoại 3D nằm ngoài phạm vi thực hiện của nghiên cứu này. Việc tập trung vào các giới hạn kỹ thuật nêu trên giúp hệ thống đảm bảo tính ổn định và hiệu suất tốt nhất trong phân khúc thiết bị nhúng IoT, thay vì cố gắng dàn trải sang các ứng dụng thương mại quy mô lớn đòi hỏi hạ tầng máy chủ phức tạp. Mục tiêu cốt lõi là làm chủ công nghệ xử lý ảnh trên chip và xây dựng luồng dữ liệu an toàn, hiệu quả từ thiết bị về ứng dụng di động cho người dùng cuối.

\section{Ý nghĩa khoa học và thực tiễn}

\subsection{Ý nghĩa khoa học}
Đề tài góp phần thực nghiệm và đánh giá định lượng năng lực tính toán tại biên của kiến trúc vi xử lý nhúng thế hệ mới. Cụ thể, nghiên cứu làm rõ hiệu năng của tập lệnh mở rộng vector trên lõi Xtensa LX7 khi xử lý các cấu trúc mạng thần kinh nhân tạo tích chập hạng nhẹ như MTMN và MobileFaceNet. Kết quả thực hiện là cơ sở tham khảo quý giá cho việc tối ưu hóa các mô hình học sâu kích thước nhỏ, chứng minh tính khả thi của xu hướng Edge AI trong việc giải quyết bài toán thị giác máy tính phức tạp mà không cần phụ thuộc vào tài nguyên điện toán đám mây.

\subsection{Ý nghĩa thực tiễn}
Về mặt thực tiễn, đề tài mở ra một giải pháp công nghệ có khả năng thương mại hóa cao với chi phí sản xuất cực kỳ tối ưu. Hệ thống khóa thông minh nhận diện khuôn mặt này hoàn toàn phù hợp để ứng dụng trực tiếp tại các mô hình phòng trọ sinh viên, chung cư mini hoặc các văn phòng doanh nghiệp nhỏ. Sự kết hợp giữa phần cứng giá thành rẻ và luồng quản lý đồng bộ qua ứng dụng di động giúp người dùng tiếp cận công nghệ an ninh cao cấp một cách dễ dàng, đồng thời giải quyết triệt để bài toán quản lý ra vào một cách an toàn, tiện lợi cho các chủ cơ sở kinh doanh.

\section{Cấu trúc của báo cáo}
Nội dung của báo cáo được chia thành 5 chương:
\begin{enumerate}
    \item \textbf{Chương 1: Giới thiệu.} Trình bày lý do chọn đề tài, mục tiêu, đối tượng, phạm vi nghiên cứu và ý nghĩa của đề tài.
    \item \textbf{Chương 2: Cơ sở lý thuyết và công nghệ sử dụng.} Cung cấp các kiến thức nền tảng về thị giác máy tính, TinyML, cùng các công cụ phần cứng và phần mềm được áp dụng.
    \item \textbf{Chương 3: Phân tích và thiết kế hệ thống.} Mô tả chi tiết yêu cầu bài toán, kiến trúc hệ thống tổng thể, thiết kế cơ sở dữ liệu và các biểu đồ mô hình hóa.
    \item \textbf{Chương 4: Triển khai thử nghiệm và đánh giá.} Trình bày quá trình thi công thực tế, kịch bản kiểm thử và các số liệu phân tích, đánh giá hiệu năng của hệ thống.
    \item \textbf{Chương 5: Kết luận và hướng phát triển.} Tổng kết những kết quả đã đạt được, nêu ra hạn chế còn tồn đọng và đề xuất các hướng phát triển nâng cấp trong tương lai.
\end{enumerate}